% ============================================================
\section{Architecture Assessment and Roadmap}
\label{sec:arch-assessment}
% ============================================================

This chapter records an external architectural review of the
\texttt{hacdcpf} package conducted against the updated codebase described
in Chapters~\ref{sec:data-model-deep}--\ref{sec:analysis}.
It documents observed strengths, remaining design risks, module-level
recommendations, testing strategy, dependency rules, and a prioritised
roadmap.

% ------------------------------------------------------------
\subsection{Overall Quality Assessment}
\label{sec:quality-table}
% ------------------------------------------------------------

The package architecture now resembles a mature scientific computing library.
Table~\ref{tab:quality} summarises the current design quality by area.

\begin{longtable}{lrl}
\caption{Architecture quality by area.\label{tab:quality}} \\
\hline\hline
\textbf{Area} & \textbf{Quality} & \textbf{Comment} \\
\hline\endfirsthead
\hline\hline
\textbf{Area} & \textbf{Quality} & \textbf{Comment} \\
\hline\endhead
\hline\endfoot
Public model design & Strong & Rich component model; clear I/O field dictionaries \\
Solver layering & Strong & Projection and assembly separation is a major improvement \\
PF solver design & Strong & Multiple routes and fallback methods well documented \\
OPF design & Good--Strong & Native, parity-IPM, DC-QP, and reactive routes well separated \\
I/O contracts & Strong & JSON and Excel field dictionaries are detailed \\
Error handling & Improved & \texttt{Result<T>}, \texttt{Error}, \texttt{ErrorCode} are good foundations \\
Analysis modules & Promising & Requires careful dependency control \\
Backend abstraction & Improved & Linear/NLP/QP backend split is the right direction \\
Documentation & Strong & Implementation-facing handbook is detailed \\
Testability & Likely improved & Needs evidence through test matrix and CI policy \\
\end{longtable}

The overall data-flow pipeline is now clearly separated:

\begin{center}
\texttt{HybridPowerSystem}
$\;\downarrow\;$
Projection Layer
$\;\downarrow\;$
Canonical Models
$\;\downarrow\;$
Assembly Layer
$\;\downarrow\;$
Flat Solver Data
$\;\downarrow\;$
PF / OPF / Analysis Solvers
$\;\downarrow\;$
Result Objects
$\;\downarrow\;$
I/O and Reports
\end{center}

% ------------------------------------------------------------
\subsection{Confirmed Strengths}
\label{sec:confirmed-strengths}
% ------------------------------------------------------------

\subsubsection{Horizontal Infrastructure Layer}

The dedicated solver-infrastructure layer covering assembly, projection,
backend abstractions, and the thread pool (Chapter~\ref{sec:solver-infra})
is a major improvement.  This separation is essential for long-term
maintainability: the rich physical model can evolve without destabilising
the numerical kernels.

\subsubsection{Monadic Result Handling}

The \texttt{Result<T>} / \texttt{Error} / \texttt{ErrorCode} design
(Section~\ref{sec:error-result}) addresses one of the most important concerns in
simulation packages: handling invalid input and solver failures predictably.
A monadic result type is much better than relying on bare exceptions,
booleans, or string messages.

\subsubsection{Backward-Compatible Option Refactoring}

\texttt{PowerFlowOptions::from\_parts()} preserves flat API compatibility for
existing users while enabling structured internal configuration.  The
recommended long-term pattern is:

\begin{verbatim}
struct PowerFlowOptions {
    // Legacy flat fields
    double tolerance;
    int max_iterations;
    bool enable_homotopy;

    static PowerFlowOptions from_parts(
        const ConvergenceOptions& convergence,
        const LinearSolverOptions& linear,
        const GlobalizationOptions& globalization,
        const PVPQSwitchingOptions& pvpq,
        const ConverterModeOptions& converter,
        const ZipLoadOptions& zip,
        const RuntimeOptions& runtime
    );

    PowerFlowOptionParts to_parts() const; // optional round-trip
};
\end{verbatim}

\subsubsection{Explicit Validation Subsystem}

A dedicated validation layer capable of catching missing slack buses,
duplicate IDs, invalid branch endpoints, inconsistent voltage limits,
unsupported OPF configurations, and islands without a source gives the
package its primary defence against badly formed inputs.  The recommended
list of checks is given in Section~\ref{sec:validation-deep}.

\subsubsection{Projection plus Assembly Separation}

\begin{itemize}
  \item \textbf{Projection}: rich physical model $\to$ canonical electrical model.
  \item \textbf{Assembly}: canonical model $\to$ numerical vectors, matrices,
        and residual functions.
\end{itemize}

This boundary allows rich devices (motors, chargers, PV arrays, flexible
loads, switches, transformers, VSC and DCDC converters) to evolve independently
of solver kernels.

\subsubsection{Centralised Defaults}

Every default is now traceable to one of: engineering convention, MATPOWER
compatibility, package-specific design, safety fallback, or numerical
stabilisation.  Avoid silent defaults for safety-critical fields; such fields
should be required rather than defaulted.

% ------------------------------------------------------------
\subsection{Remaining Design Risks}
\label{sec:design-risks}
% ------------------------------------------------------------

\subsubsection{Analysis Layer May Become Too Broad}

The \texttt{analysis/} module currently covers carbon tracing, distribution
PF, three-phase unbalanced PF, PV modelling, reliability, resilience, and
time-series simulation.  These are not all the same type of functionality.
Consider explicit subdomains:

\begin{verbatim}
include/hacdcpf/analysis/
  carbon/
  distribution/
  three_phase/
  pv/
  reliability/
  resilience/
  time_series/
\end{verbatim}

\paragraph{Dependency rule for analysis}
Analysis modules may depend on:
\begin{itemize}
  \item public model and validation API,
  \item public solver API and result types.
\end{itemize}
They must not depend on private Newton/IPM internals, assembly private data,
or backend private classes.

\subsubsection{Three-Phase and Balanced PF Should Not Share One Abstraction}

Three-phase unbalanced PF is structurally different from balanced
positive-sequence PF.  Prefer distinct free functions over one overly
generic virtual interface:

\begin{verbatim}
BalancedPowerFlowResult solve_power_flow(
    const HybridPowerSystem& sys,
    const PowerFlowOptions& opt
);

ThreePhasePowerFlowResult solve_three_phase_power_flow(
    const ThreePhaseSystem& sys,
    const ThreePhasePowerFlowOptions& opt
);
\end{verbatim}

A common \texttt{SolverSummary} (converged, iterations, max residual, status)
can be shared where natural without forcing identical result shapes.

\subsubsection{Error Handling Policy Needs Explicit Definition}

Now that both \texttt{Result<T>} and exceptions are present, define exactly
when each mechanism is used.  A recommended policy is given in
Table~\ref{tab:error-policy}.

\begin{longtable}{ll}
\caption{Recommended error handling policy.\label{tab:error-policy}} \\
\hline\hline
\textbf{Situation} & \textbf{Mechanism} \\
\hline\endfirsthead\hline\hline\textbf{Situation}&\textbf{Mechanism}\\\hline\endhead\hline\endfoot
Expected input validation failure & \texttt{Result<T>} \\
Solver non-convergence & result object with status \\
Optional feature unavailable & \texttt{Result<T>} consistently \\
Programmer error / invariant violation & assertion or internal exception \\
File not found & \texttt{Result<T>} for library API \\
JSON parse failure & \texttt{Result<T>} with error code and location \\
Backend crash / failure & \texttt{Result<T>} with backend diagnostic \\
\end{longtable}

Provide both styles where helpful:
\begin{verbatim}
Result<HybridPowerSystem> try_load_json(const std::string& path);
HybridPowerSystem load_json_or_throw(const std::string& path);
\end{verbatim}

\subsubsection{Field Dictionary / Serializer Drift}

Detailed I/O field dictionaries are excellent, but documentation, serializers,
and model structs can drift apart over time.  Automated tests should verify:
\begin{itemize}
  \item every documented field exists in the serializer,
  \item every serializer field is documented,
  \item JSON round-trip preserves each field,
  \item each documented default equals the code default.
\end{itemize}

A single machine-readable field-descriptor registry
(see Section~\ref{sec:field-registry}) is the preferred long-term solution.

\subsubsection{Backend Capability Discovery}

Since the package has optional backends (HiGHS, Gurobi, Ipopt, native IPM,
etc.), expose a runtime capability query:

\begin{verbatim}
struct SolverCapabilities {
    bool has_sparse_lu;
    bool has_highs;
    bool has_gurobi;
    bool has_native_ipm;
    bool has_excel;
    bool supports_quadratic_objective;
    bool supports_integer_variables;
};

SolverCapabilities get_solver_capabilities();
\end{verbatim}

This enables portable user code that selects a backend based on availability
rather than compile-time flags alone.

% ------------------------------------------------------------
\subsection{Module-Level Recommendations}
\label{sec:module-recommendations}
% ------------------------------------------------------------

\subsubsection{Model Layer}

\paragraph{Separate persistent fields from runtime fields}
If any model struct contains both user-input data and runtime caches, split
them.  Do not mix result fields into input model objects unless explicitly
supporting in-place result annotation.

\paragraph{Strong IDs at boundaries}
At projection boundaries, typed reference structs avoid accidentally
connecting a DC converter terminal to an AC bus field:
\begin{verbatim}
struct ACBusId { std::string value; };
struct DCBusId { std::string value; };
\end{verbatim}

\paragraph{Stable enum serialisation}
Treat enum string names as serialised API.  Document whether strings are
case-sensitive.  Avoid casual renaming once a version ships.

\subsubsection{Validation Layer}

Use tiered validation levels:

\begin{verbatim}
enum class ValidationLevel {
    Basic,        // required fields, references, limits
    Electrical,   // topology, slack/reference, physical consistency
    SolverReady,  // checks required by PF/OPF kernels
    Strict        // warnings promoted to errors
};
\end{verbatim}

A structured diagnostic shape supports GUIs, APIs, and batch workflows:

\begin{verbatim}
struct Diagnostic {
    Severity severity;   // Info, Warning, Error
    ErrorCode code;
    std::string message;
    std::string component_type;
    std::string component_id;
    std::string field;
};
\end{verbatim}

\subsubsection{Projection Layer}

\paragraph{Projection report}
Every projection should optionally return a diagnostic report that records
mappings such as:
\begin{center}
\begin{tabular}{ll}
\hline\hline
\textbf{Source} & \textbf{Canonical} \\
\hline
\texttt{PVSystem pv\_12} & \texttt{StaticGenerator sg\_pv\_12} \\
\texttt{Motor motor\_7} & \texttt{EquivalentLoad load\_motor\_7} \\
\texttt{ClosedSwitch sw\_3} & \texttt{BusMerge group\_2} \\
\texttt{OpenSwitch sw\_4} & \texttt{DisabledBranch branch\_sw\_4} \\
\hline\hline
\end{tabular}
\end{center}

\paragraph{Reversible mapping}
Maintain a \texttt{ComponentMapping} record with source type, source ID,
canonical type, canonical ID, and participation factor.  This is essential
for result attribution, carbon tracing, and reporting.

\subsubsection{Assembly Layer}

Assembly should be \emph{deterministic}: given the same canonical system,
it must produce the same ordering, index maps, vector layout, and sparsity
pattern.  Explicit layout objects make residual code readable:

\begin{verbatim}
struct PowerFlowVariableLayout {
    std::vector<ACBusIndex>  angle_variables;
    std::vector<ACBusIndex>  voltage_variables;
    std::vector<DCBusIndex>  dc_voltage_variables;
    std::vector<VSCIndex>    converter_variables;
};
\end{verbatim}

\subsubsection{Power-Flow Layer}

\paragraph{Standardise solver status}
Every PF route should return a common summary:
\begin{verbatim}
struct SolverSummary {
    TerminationStatus status;
    bool   converged;
    int    iterations;
    double max_mismatch;
    double elapsed_ms;
    std::vector<Diagnostic> diagnostics;
};
\end{verbatim}

\paragraph{Visible fallback trace}
Record every solver method transition in the result:
\begin{verbatim}
{
  "solver_trace": [
    {"method":"Newton",
     "status":"SingularJacobian", "iterations":4},
    {"method":"HomotopyContinuation",
     "status":"Converged", "iterations":12}
  ]
}
\end{verbatim}

\subsubsection{OPF Layer}

\paragraph{Formulation separate from solver}
For each OPF route maintain the pipeline:
model extraction $\to$ formulation construction $\to$ backend solve
$\to$ result mapping $\to$ post-check.

\paragraph{Independent feasibility audit}
After every solve, independently check AC/DC power balance, voltage limits,
generator and branch limits, converter limits, and objective recomputation.
Return a structured \texttt{OpfAudit} object.

\paragraph{Explainable infeasibility}
Provide likely-cause diagnostics instead of a bare ``infeasible'' status:
\begin{itemize}
  \item bus has load but no connected source,
  \item generator \texttt{p\_min\_mw} $>$ \texttt{p\_max\_mw},
  \item voltage upper limit below lower limit,
  \item converter power exceeds DC grid capacity.
\end{itemize}

\subsubsection{I/O Layer}
\label{sec:io-recommendations}

\paragraph{Format versioning}
Every exported JSON should include:
\begin{verbatim}
{
  "format": "hacdcpf.hybrid_power_system",
  "format_version": "1.0",
  "package_version": "x.y.z"
}
\end{verbatim}

\paragraph{Import modes}
\begin{verbatim}
enum class ImportMode { Strict, Permissive };
\end{verbatim}
In Strict mode unknown columns are errors; in Permissive mode they become
warnings.

\paragraph{Field-descriptor registry}
\label{sec:field-registry}

A machine-readable registry derives the serialiser, documentation table, and
tests from one source:
\begin{verbatim}
constexpr FieldDescriptor ac_bus_fields[] = {
    {"id",    Required::Yes, DefaultPolicy::None},
    {"vn_kv", Required::Yes, DefaultPolicy::None},
    {"vm_pu", Required::No,  DefaultPolicy::One},
};
\end{verbatim}

\subsubsection{Analysis Suite}

\paragraph{Carbon tracing}
Expose the tracing method explicitly:
\begin{verbatim}
enum class CarbonTracingMethod {
    AverageMix, ProportionalSharing,
    MarginalEmission, FlowTracing
};
\end{verbatim}

\paragraph{Time-series simulation}
Avoid mutating the base system in uncontrolled ways.  Recommended pattern:
\begin{verbatim}
struct TimeSeriesScenario {
    HybridPowerSystem base_system;
    std::vector<TimeStepInput> steps;
};
\end{verbatim}

% ------------------------------------------------------------
\subsection{Recommended Testing Matrix}
\label{sec:test-matrix}
% ------------------------------------------------------------

Testing should be organised by layer.
Table~\ref{tab:test-matrix} lists the recommended test groups.

\begin{longtable}{ll}
\caption{Recommended test groups by layer.\label{tab:test-matrix}} \\
\hline\hline
\textbf{Layer} & \textbf{Test groups} \\
\hline\endfirsthead\hline\hline\textbf{Layer}&\textbf{Test groups}\\\hline\endhead\hline\endfoot
Model / validation & ModelDefaultsTest, EnumSerializationTest, \\
 & ValidationDuplicateIdsTest, ValidationDanglingReferencesTest, \\
 & ValidationInvalidLimitsTest, ValidationIslandWithoutSlackTest \\
Projection & ProjectionSwitchMergingTest, ProjectionMotorAsLoadTest, \\
 & ProjectionPVAsStaticGeneratorTest, ProjectionStorageMappingTest \\
Assembly & AssemblyVariableOrderingTest, AssemblySparsePatternTest, \\
 & AssemblyIndexMapTest, AssemblyPerUnitConversionTest \\
Power flow & PFIEEE14ConvergenceTest, PFHybridTwoBusVSCTest, \\
 & PFDCGridConvergenceTest, PFIslandedSystemTest, \\
 & PFHomotopyFallbackTest, PFNewtonKrylovFallbackTest, \\
 & PFNCPComplementarityTest \\
OPF & OPFACSmallCaseTest, OPFDCQPBackendTest, OPFParityIPMTest, \\
 & OPFReactiveTapShuntTest, OPFInfeasibilityDiagnosticTest \\
I/O & JsonRoundTripAllComponentsTest, MatpowerImportIEEE14Test, \\
 & UnknownFieldStrictModeTest, MissingOptionalFieldDefaultTest, \\
 & ResultJsonSchemaTest \\
Analysis & CarbonTracingKnownCaseTest, DistributionRadialFeederTest, \\
 & ThreePhaseUnbalancedKnownCaseTest, PVArrayIrradianceCurveTest, \\
 & ReliabilitySingleOutageTest, TimeSeriesProfileReplayTest \\
\end{longtable}

For complementarity variables $a$ and $b$, each NCP unit test should verify:
\begin{equation}
  \phi(a,b)=0 \iff a\ge 0,\; b\ge 0,\; ab=0.
\end{equation}

% ------------------------------------------------------------
\subsection{Recommended CI Pipeline}
\label{sec:ci-pipeline}
% ------------------------------------------------------------

\begin{center}
\begin{tabular}{cl}
\hline\hline
\textbf{Stage} & \textbf{Jobs} \\
\hline
1 & \texttt{clang-format} check \\
2 & \texttt{clang-tidy} selected checks \\
3 & Build matrix: Debug / Release, Excel OFF / ON, backends disabled / enabled \\
4 & Unit tests: model, validation, projection, assembly, I/O \\
5 & Numerical tests: PF, OPF, analysis \\
6 & Sanitizers: AddressSanitizer, UndefinedBehaviorSanitizer \\
7 & Documentation: LaTeX build \\
\hline
Nightly & Large benchmarks, ThreadSanitizer, coverage, performance regression \\
\hline\hline
\end{tabular}
\end{center}

% ------------------------------------------------------------
\subsection{Dependency Boundary Rules}
\label{sec:dep-rules}
% ------------------------------------------------------------

\subsubsection{Allowed Dependencies}

\begin{center}
\begin{tabular}{ll}
\hline\hline
\textbf{Module} & \textbf{May depend on} \\
\hline
\texttt{api} & \texttt{model, io, validation, power\_flow, optimal\_power\_flow, analysis} \\
\texttt{io} & \texttt{model, validation, result types} \\
\texttt{validation} & \texttt{model} \\
\texttt{projection} & \texttt{model, validation} \\
\texttt{assembly} & canonical/projection output, model enums, utilities \\
\texttt{power\_flow} & \texttt{assembly, backend, model result types} \\
\texttt{optimal\_power\_flow} & \texttt{assembly, backend, power\_models} \\
\texttt{power\_models} & AML, canonical model data \\
\texttt{analysis} & \texttt{model, result types, public solver APIs} \\
\texttt{backend} & external solver libraries \\
\texttt{engine} & compatibility/adapters only \\
\hline\hline
\end{tabular}
\end{center}

\subsubsection{Disallowed Dependencies}

\begin{center}
\begin{tabular}{ll}
\hline\hline
\textbf{Module} & \textbf{Must not depend on} \\
\hline
\texttt{model} & \texttt{io, power\_flow, optimal\_power\_flow, analysis} \\
\texttt{validation} & solver internals \\
\texttt{power\_flow} & \texttt{io} \\
\texttt{optimal\_power\_flow} & \texttt{io} \\
\texttt{assembly} & Excel / JSON I/O \\
\texttt{backend} & model-specific business logic \\
\texttt{analysis} & private Newton/IPM internals (unless intentionally internal) \\
\hline\hline
\end{tabular}
\end{center}

These boundaries can be enforced by include-guard audits or a build-system
layer-visibility policy.

% ------------------------------------------------------------
\subsection{Priority Roadmap}
\label{sec:roadmap}
% ------------------------------------------------------------

\paragraph{Priority 1 --- Trustworthiness}
\begin{enumerate}
  \item Add independent PF and OPF audit checks (\texttt{OpfAudit}).
  \item Add derivative finite-difference tests for Jacobians.
  \item Add serialisation completeness tests.
  \item Add projection mapping reports.
  \item Add solver trace reporting for fallback method transitions.
\end{enumerate}

\paragraph{Priority 2 --- User Experience}
\begin{enumerate}
  \item Add strict/permissive import modes.
  \item Add runtime capability discovery API (\texttt{get\_solver\_capabilities()}).
  \item Add better infeasibility diagnostics with probable causes.
  \item Add example workflows (see Section~\ref{sec:example-workflow}).
  \item Publish JSON schema files.
\end{enumerate}

\paragraph{Priority 3 --- Maintainability}
\begin{enumerate}
  \item Enforce dependency boundaries via build-system rules.
  \item Split large analysis submodules if necessary.
  \item Centralise field descriptors into a machine-readable registry.
  \item Keep \texttt{engine/} as adapter/fa\c{c}ade only.
\end{enumerate}

\paragraph{Priority 4 --- Performance and Scalability}
\begin{enumerate}
  \item Benchmark assembly time separately from solve time.
  \item Track sparse-matrix nonzero counts across test cases.
  \item Add performance regression cases for large networks.
  \item Profile three-phase and time-series simulation.
  \item Cache symbolic factorisations where safe.
\end{enumerate}

% ------------------------------------------------------------
\subsection{Example Clean Public Workflow}
\label{sec:example-workflow}
% ------------------------------------------------------------

A mature public API should look like the following.  The pattern
``load, validate, configure, solve, report'' captures the full user
journey:

\begin{verbatim}
#include <hacdcpf/api/hacdcpf.hpp>

int main() {
    // 1. Load with error propagation
    auto system_result = hacdcpf::try_load_json("case.json");
    if (!system_result) {
        std::cerr << system_result.error().message << "\n";
        return 1;
    }
    auto system = system_result.value();

    // 2. Validate before solving
    auto report = hacdcpf::validate(
        system, hacdcpf::ValidationLevel::SolverReady);
    if (!report.ok()) {
        std::cerr << report.summary() << "\n";
        return 1;
    }

    // 3. Structured option assembly
    auto options = hacdcpf::PowerFlowOptions::from_parts(
        hacdcpf::ConvergenceOptions{},
        hacdcpf::LinearSolverOptions{},
        hacdcpf::GlobalizationOptions{},
        hacdcpf::PVPQSwitchingOptions{},
        hacdcpf::ConverterModeOptions{},
        hacdcpf::ZipLoadOptions{},
        hacdcpf::RuntimeOptions{}
    );

    // 4. Solve
    auto pf = hacdcpf::solve_power_flow(system, options);
    if (!pf.converged) {
        std::cerr << pf.diagnostics.termination_reason << "\n";
        return 2;
    }

    // 5. Save result
    hacdcpf::save_power_flow_result_json(pf, "pf_result.json");

    // 6. Post-solve analysis
    auto carbon = hacdcpf::analysis::run_carbon_analysis(system, options);
    hacdcpf::analysis::save_carbon_report_json(carbon, "carbon.json");

    return 0;
}
\end{verbatim}

% ------------------------------------------------------------
\subsection{Final Recommendation Summary}
\label{sec:final-recommendation}
% ------------------------------------------------------------

The package has moved substantially in the right direction.
The next most valuable work is not adding more solver features,
but improving \textbf{trust, diagnostics, and maintainability}:

\begin{enumerate}
  \item Add independent numerical audit reports for PF and OPF.
  \item Add automated serialisation completeness tests.
  \item Add projection mapping reports.
  \item Add runtime backend capability discovery.
  \item Enforce dependency boundary rules.
  \item Keep analysis modules modular and subdomain-separated.
  \item Improve infeasibility and non-convergence explanations.
  \item Build benchmark and regression infrastructure.
\end{enumerate}

The main remaining structural risks are:

\begin{itemize}
  \item \texttt{analysis/} becoming a miscellaneous collection without
        subdomain separation.
  \item Documentation and serialisers drifting without a field-descriptor
        registry.
  \item Incomplete backend capability discovery forcing compile-time
        conditional code into user applications.
  \item Three-phase unbalanced models being over-coupled to balanced PF
        abstractions.
  \item OPF infeasibility diagnostics remaining too coarse for user debugging.
\end{itemize}

If the roadmap in Section~\ref{sec:roadmap} is followed, the package will be
significantly easier to maintain, debug, validate, and adopt by external users.
